
Neural networks trained with empirical risk minimization on group-imbalanced data exhibit poor worst-group accuracy. Prior work established that minority groups occupy sharper regions of the loss landscape (sharpness ratio SR > 1.0), but the origin of this curvature asymmetry remained unclear. This work investigates the Differential Convergence Rate hypothesis: majority groups converge faster during training, smoothing curvature in majority-relevant parameter directions while leaving minority directions sharp.

Two findings are presented. First, sharpness ratio equals approximately 1.0 at random initialization ($SR_0$ = 0.9999, 95\% CI [0.9953, 1.0046], n = 5 seeds), indicating that intrinsic data geometry does not cause curvature asymmetry. Second, gradient ratio decay temporally precedes sharpness ratio divergence by $\tau$ = 3.67 epochs (95\% CI [3.0, 4.0], n = 3 seeds). This positive lag is consistent with a causal mechanism wherein majority convergence precedes minority curvature increase.

These findings characterize a diagnostic window: gradient dynamics signal impending curvature divergence 3-4 epochs before it manifests. A proposed update-norm parity intervention was implemented and code-validated but not experimentally executed due to hardware constraints.

